% Frontier relation-map: rows 1-2 of the old fig-r4-relation.tex, split out.
%
% SPLIT (noted, not silently): the single R4 relation-map carried three rows
% -- the rate-distortion constraint upgrade, the delta=0 entropy corner, and
% the twenty-questions/zoom correspondence -- while its caption described
% only the third, and it was cited from Sec. sec:zoom mid-analogy. A reader
% arriving there was handed a diagram whose top two-thirds were about the
% previous subsection. Rows 1-2 now live here, cited from sec:rdf where they
% belong; row 3 stays in fig-r4-relation.tex, cited from sec:zoom. Same
% grammar, same unit length, same content -- only the placement changed.
% [internal, script r4_figures.py]
\begin{figure}[t]
\centering
\begin{tikzpicture}[x=1cm,y=1cm]
\node[align=center,font=\small\bfseries,text=harborblue,text width=12cm] at (6,8.25)
  {R4a --- the operator's dial: a second constraint buys back bits, and $\delta{=}0$ gives them all back};

% Base domain (left): classical single-constraint rate-distortion
\draw[draw=harborblue,thick,fill=harborblue,fill opacity=0.12] (0.1,1.25) rectangle (5.0,7.65);
\node[align=center,font=\small\bfseries,text width=4.5cm] at (2.55,7.2) {Base: single-constraint RD};

% Target domain (right): the two-constraint dial
\draw[draw=seagreen,thick,fill=seagreen,fill opacity=0.12] (7.0,1.25) rectangle (11.9,7.65);
\node[align=center,font=\small\bfseries,text width=4.5cm] at (9.45,7.2) {Target: $R(\delta,f)$, the operator dial};

% Row 1
\node[align=center,font=\scriptsize,text width=4.3cm] at (2.55,5.3) {minimize $I(X;\hat X)$\\s.t.\ distortion $\le D$\\(Shannon 1959, textbook RDC)};
\node[align=center,font=\scriptsize,text width=4.3cm] at (9.45,5.3) {minimize $I(X;\hat X)$\\s.t.\ miss rate $\le \delta$\\AND flag rate $\le f$\\(custom formulation)};
\draw[<->,>=Stealth,thick,shipred] (5.15,5.0) -- (6.85,5.0);
\node[align=center,font=\scriptsize\bfseries,text=shipred,text width=1.85cm] at (6.0,6.15) {a SECOND knob:\\flags buy back bits};

% Row 2
\node[align=center,font=\scriptsize,text width=4.3cm] at (2.55,2.9) {zero tolerance:\\$\delta = 0$, no missed\\positives allowed};
\node[align=center,font=\scriptsize,text width=4.3cm] at (9.45,2.9) {$R(0, f{\to}p) = H(p)$\\$p{=}0.05 \Rightarrow 0.286$ bits/sym\\{[verified, Cover--Thomas]}\\the strict corner, and Theorem 1's own corner};
\draw[<->,>=Stealth,thick,shipred] (5.15,2.6) -- (6.85,2.6);
\node[align=center,font=\scriptsize\bfseries,text=shipred,text width=1.85cm] at (6.0,3.7) {$\delta{=}0$ boundary\\$\Leftrightarrow$ entropy floor};

\node[align=center,font=\scriptsize,text width=11.5cm] at (6.0,0.35)
  {\itshape Two faces of one dial: the constraint that got added, and the boundary where it collapses back to Shannon. $R(0,0.10)=0.186$, $R(0.01,0.10)=0.110$, $R(0.04,0.06)=0.009$ bits/sym [internal, \texttt{b1\_frontier.py}].};
\end{tikzpicture}
\caption{The relation-map for the two-constraint program. The relation that transfers from classical rate--distortion is not the objective --- that is unchanged, $I(X;\hat X)$ --- but the shape of the feasible set: pricing false negatives separately from flag rate adds a second binding constraint, and every bit the digest saves is bought by loosening one of the two. Row 2 is the seam with Theorem 1: at $\delta{=}0$ and $f\to p$ the program returns the source's entire entropy, which is the frontier's own version of the covering floor.}
\label{fig:rdfrel}
\end{figure}
